% GENERATED FILE. Edit canonical lessons, then run npm run generate:books.
% canonical-source-hash: fnv1a64:bf8ca717ea8cdc34
% canonical-lessons: KA-C531-pagade, KA-C531-keram, KA-C531-kannamuccale, KA-C531-uyyale, KA-C531-atike

\chapter{Games at Home}
\label{ch:ka-games-at-home}
\hlchaptermodality{531}

\begin{chapteropening}
\textbf{By the end of this chapter:} \emph{I can say a dice game, carrom, hide-and-seek, a swing and a toy.}

Complete the last lesson of chapter 531: I can say a dice game, carrom, hide-and-seek, a swing and a toy.
\end{chapteropening}

\section[\texorpdfstring{\kn{ಪಗಡೆ} (pagaḍe)}{ಪಗಡೆ (pagaḍe)}]{\kn{ಪಗಡೆ} (pagaḍe) — a dice game}

\label{lesson:KA-C531-pagade}



\begin{quote}
Before the new one: say the Kannada for to savour, then the Kannada for to be absorbed.
\end{quote}

\subsection*{You'll want to know: \kn{ಪಗಡೆ}}
\textbf{\kn{ಪಗಡೆ}} — \emph{pagaḍe} — \textquotedblleft{}a dice game\textquotedblright{}.

\begin{quote}
\kn{ಹಬ್ಬದ} \kn{ದಿನ} \kn{ಅಜ್ಜಿ} \kn{ಮತ್ತು} \kn{ಅತ್ತೆ} \kn{ಪಗಡೆ} \kn{ಆಡಿದರು}.
\end{quote}

\begin{quote}
\emph{habbada dina ajji mattu atte pagaḍe āḍidaru.}
\end{quote}

\begin{quote}
On the festival day grandmother and aunt played pagade.
\end{quote}

A board game played with cowries or long dice.

\begin{cousinweb}
\textbf{\kn{ಪಗಡೆ}} (\emph{pagaḍe}) is the name of an old board game played with dice; DEDR 1461 records the game, but the origin of the name is not settled.
\end{cousinweb}

\begin{grammarlens}[title={What you've built}]
One more word for describing what you lived and felt.
\end{grammarlens}

\subsection*{Your turn}
\begin{itemize}
\raggedright
  \item \emph{Say it:} \emph{pagaḍe}
  \item \emph{Say it:} \emph{pagaḍe}, once more
  \item \emph{Say it:} \emph{habbada dina ajji mattu atte pagaḍe āḍidaru.} — On the festival day grandmother and aunt played pagade.
  \item \emph{Recall:} say the Kannada for to savour, then the Kannada for to be absorbed, then say \emph{pagaḍe} again
\end{itemize}

\begin{tcolorbox}[breakable,colback=teal!4,colframe=teal!35!black,title={Before you move on}]
What does \kn{ಪಗಡೆ} mean? (\textquotedblleft{}A dice game\textquotedblright{}.) Say it once more.
\end{tcolorbox}
\section[\texorpdfstring{\kn{ಕೇರಂ} (kēraṃ)}{ಕೇರಂ (kēraṃ)}]{\kn{ಕೇರಂ} (kēraṃ) — carrom}

\label{lesson:KA-C531-keram}



\begin{quote}
Before the new one: say the Kannada for to be absorbed, then the Kannada for a dice game.
\end{quote}

\subsection*{You'll want to know: \kn{ಕೇರಂ}}
\textbf{\kn{ಕೇರಂ}} — \emph{kēraṃ} — \textquotedblleft{}carrom\textquotedblright{}.

\begin{quote}
\kn{ರಜೆಯಲ್ಲಿ} \kn{ಮಕ್ಕಳು} \kn{ದಿನವೂ} \kn{ಕೇರಂ} \kn{ಆಡುತ್ತಾರೆ}.
\end{quote}

\begin{quote}
\emph{rajeyalli makkaḷu dinavū kēraṃ āḍuttāre.}
\end{quote}

\begin{quote}
In the holidays the children play carrom every day.
\end{quote}

Players flick a striker to pocket the pieces.

\begin{cousinweb}
\textbf{\kn{ಕೇರಂ}} (\emph{kēraṃ}) is a loan from English \emph{carrom}, the name of the board game.
\end{cousinweb}

\begin{grammarlens}[title={What you've built}]
One more word for describing what you lived and felt.
\end{grammarlens}

\subsection*{Your turn}
\begin{itemize}
\raggedright
  \item \emph{Say it:} \emph{kēraṃ}
  \item \emph{Say it:} \emph{kēraṃ}, once more
  \item \emph{Say it:} \emph{rajeyalli makkaḷu dinavū kēraṃ āḍuttāre.} — In the holidays the children play carrom every day.
  \item \emph{Recall:} say the Kannada for to be absorbed, then the Kannada for a dice game, then say \emph{kēraṃ} again
\end{itemize}

\begin{tcolorbox}[breakable,colback=teal!4,colframe=teal!35!black,title={Before you move on}]
What does \kn{ಕೇರಂ} mean? (\textquotedblleft{}Carrom\textquotedblright{}.) Say it once more.
\end{tcolorbox}
\section[\texorpdfstring{\kn{ಕಣ್ಣಾಮುಚ್ಚಾಲೆ} (kaṇṇāmuccāle)}{ಕಣ್ಣಾಮುಚ್ಚಾಲೆ (kaṇṇāmuccāle)}]{\kn{ಕಣ್ಣಾಮುಚ್ಚಾಲೆ} (kaṇṇāmuccāle) — hide-and-seek}

\label{lesson:KA-C531-kannamuccale}



\begin{quote}
Before the new one: say the Kannada for a dice game, then the Kannada for carrom.
\end{quote}

\subsection*{You'll want to know: \kn{ಕಣ್ಣಾಮುಚ್ಚಾಲೆ}}
\textbf{\kn{ಕಣ್ಣಾಮುಚ್ಚಾಲೆ}} — \emph{kaṇṇāmuccāle} — \textquotedblleft{}hide-and-seek\textquotedblright{}.

\begin{quote}
\kn{ಮಕ್ಕಳು} \kn{ತೋಟದಲ್ಲಿ} \kn{ಕಣ್ಣಾಮುಚ್ಚಾಲೆ} \kn{ಆಡುತ್ತಿದ್ದರು}.
\end{quote}

\begin{quote}
\emph{makkaḷu tōṭadalli kaṇṇāmuccāle āḍuttiddaru.}
\end{quote}

\begin{quote}
The children were playing hide-and-seek in the garden.
\end{quote}

From \kn{ಕಣ್ಣು}, eye, and \kn{ಮುಚ್ಚು}, to close: the seeker closes his eyes.

\begin{cousinweb}
\textbf{\kn{ಕಣ್ಣಾಮುಚ್ಚಾಲೆ}} (\emph{kaṇṇāmuccāle}) is a native Kannada game name built on \kn{ಕಣ್ಣು}, \textquotedblleft{}eye\textquotedblright{}, and \kn{ಮುಚ್ಚು}, \textquotedblleft{}to close\textquotedblright{}; the origin of the ending is not settled.
\end{cousinweb}

\begin{grammarlens}[title={What you've built}]
One more word for describing what you lived and felt.
\end{grammarlens}

\subsection*{Your turn}
\begin{itemize}
\raggedright
  \item \emph{Say it:} \emph{kaṇṇāmuccāle}
  \item \emph{Say it:} \emph{kaṇṇāmuccāle}, once more
  \item \emph{Say it:} \emph{makkaḷu tōṭadalli kaṇṇāmuccāle āḍuttiddaru.} — The children were playing hide-and-seek in the garden.
  \item \emph{Recall:} say the Kannada for a dice game, then the Kannada for carrom, then say \emph{kaṇṇāmuccāle} again
\end{itemize}

\begin{tcolorbox}[breakable,colback=teal!4,colframe=teal!35!black,title={Before you move on}]
What does \kn{ಕಣ್ಣಾಮುಚ್ಚಾಲೆ} mean? (\textquotedblleft{}Hide-and-seek\textquotedblright{}.) Say it once more.
\end{tcolorbox}
\section[\texorpdfstring{\kn{ಉಯ್ಯಾಲೆ} (uyyāle)}{ಉಯ್ಯಾಲೆ (uyyāle)}]{\kn{ಉಯ್ಯಾಲೆ} (uyyāle) — a swing}

\label{lesson:KA-C531-uyyale}



\begin{quote}
Before the new one: say the Kannada for carrom, then the Kannada for hide-and-seek.
\end{quote}

\subsection*{You'll want to know: \kn{ಉಯ್ಯಾಲೆ}}
\textbf{\kn{ಉಯ್ಯಾಲೆ}} — \emph{uyyāle} — \textquotedblleft{}a swing\textquotedblright{}.

\begin{quote}
\kn{ಉದ್ಯಾನದಲ್ಲಿ} \kn{ಮಕ್ಕಳು} \kn{ಉಯ್ಯಾಲೆ} \kn{ಆಡಿದರು}.
\end{quote}

\begin{quote}
\emph{udyānadalli makkaḷu uyyāle āḍidaru.}
\end{quote}

\begin{quote}
The children played on the swings in the park.
\end{quote}

\kn{ಉಯ್ಯಾಲೆ} \kn{ಆಡು}, to swing.

\begin{cousinweb}
\textbf{\kn{ಉಯ್ಯಾಲೆ}} (\emph{uyyāle}) is native Dravidian: DEDR 731 lists Kannada \emph{uyyal(u)}, \emph{uyyāle}, \textquotedblleft{}a swing\textquotedblright{}, beside Tamil \emph{ūñcal} and Telugu \emph{uyyāla}.
\end{cousinweb}

\begin{grammarlens}[title={What you've built}]
One more word for describing what you lived and felt.
\end{grammarlens}

\subsection*{Your turn}
\begin{itemize}
\raggedright
  \item \emph{Say it:} \emph{uyyāle}
  \item \emph{Say it:} \emph{uyyāle}, once more
  \item \emph{Say it:} \emph{udyānadalli makkaḷu uyyāle āḍidaru.} — The children played on the swings in the park.
  \item \emph{Recall:} say the Kannada for carrom, then the Kannada for hide-and-seek, then say \emph{uyyāle} again
\end{itemize}

\begin{tcolorbox}[breakable,colback=teal!4,colframe=teal!35!black,title={Before you move on}]
What does \kn{ಉಯ್ಯಾಲೆ} mean? (\textquotedblleft{}A swing\textquotedblright{}.) Say it once more.
\end{tcolorbox}
\section[\texorpdfstring{\kn{ಆಟಿಕೆ} (āṭike)}{ಆಟಿಕೆ (āṭike)}]{\kn{ಆಟಿಕೆ} (āṭike) — a toy}

\label{lesson:KA-C531-atike}



\begin{quote}
Before the new one: say the Kannada for hide-and-seek, then the Kannada for a swing.
\end{quote}

\subsection*{You'll want to know: \kn{ಆಟಿಕೆ}}
\textbf{\kn{ಆಟಿಕೆ}} — \emph{āṭike} — \textquotedblleft{}a toy\textquotedblright{}.

\begin{quote}
\kn{ಅಜ್ಜ} \kn{ಮೊಮ್ಮಗನಿಗೆ} \kn{ಮರದ} \kn{ಆಟಿಕೆ} \kn{ತಂದರು}.
\end{quote}

\begin{quote}
\emph{ajja mommaganige marada āṭike tandaru.}
\end{quote}

\begin{quote}
Grandfather brought his grandson a wooden toy.
\end{quote}

Plural: \kn{ಆಟಿಕೆಗಳು}, \emph{āṭikegaḷu}.

Built on \kn{ಆಟ}, play.

\begin{cousinweb}
\textbf{\kn{ಆಟಿಕೆ}} (\emph{āṭike}) is native Dravidian: DEDR 347 lists Kannada \kn{ಆಡು}, \textquotedblleft{}to move, to play\textquotedblright{}, with the nouns \emph{āṭa} and \emph{āṭike}, \textquotedblleft{}motion, play\textquotedblright{}; a toy is the thing you play with.
\end{cousinweb}

\begin{grammarlens}[title={What you've built}]
One more word for describing what you lived and felt.
\end{grammarlens}

\subsection*{Your turn}
\begin{itemize}
\raggedright
  \item \emph{Say it:} \emph{āṭike}
  \item \emph{Say it:} \emph{āṭike}, once more
  \item \emph{Say it:} \emph{ajja mommaganige marada āṭike tandaru.} — Grandfather brought his grandson a wooden toy.
  \item \emph{Recall:} say the Kannada for hide-and-seek, then the Kannada for a swing, then say \emph{āṭike} again
\end{itemize}

\begin{tcolorbox}[breakable,colback=teal!4,colframe=teal!35!black,title={Before you move on}]
What does \kn{ಆಟಿಕೆ} mean? (\textquotedblleft{}A toy\textquotedblright{}.) Say it once more.
\end{tcolorbox}
